\BourbakiExerciseGroup

\begin{enumerate}
\def\labelenumi{\arabic{enumi})}
\item
试证：\(\mathbf{R}\) 的每个有界半开区间同胚于 \(\mathbf{R}\) 的每个无界闭区间。若 \(a < b\) ，则三个区间 \(]a, b[, [a, b[, [a, b]]\) 中任何两个都互不同胚。
\item
试证：映射 \(x \to x / (1 + |x|)\) 是 \(\mathbf{R}\) 到开区间 \(] - 1, 1[\) 上的一个同胚。
\item
设 I 表示区间 \([0, 1]\) （ \(\mathbf{R}\) 中的），并设 \(f\) 是 I 到其自身的一个同胚，满足 \(f(0) = 0\) 与 \(f(1) = 1\) 。试证：存在一个连续映射 \(g\) ，它把 \(\mathbf{I} \times \mathbf{I}\) 映入 I，使得：(i) 对每个 \(x \in \mathbf{I}\) ， \(g(x, 0) = x\) 且 \(g(x, 1) = f(x)\) ；(ii) 对每个 \(y \in \mathbf{I}\) ，偏映射 \(x \to g(x, y)\) 是 I 到其自身的同胚，满足 \(g(0, y) = 0\) 与 \(g(1, y) = 1\) 。
\item
试证：\(\mathbf{R}\) 上由 \(\overline{\mathbf{R}}\) 的唯一一致结构诱导的一致结构严格地粗于 \(\mathbf{R}\) 的加法一致结构。
\item
当点 \((x, y)\) 趋于 \((+\infty, -\infty)\) （同时保持在 \(\mathbf{R} \times \mathbf{R}\) 中）时，试证函数 \(x + y\) 的聚点所成的集是 \(\overline{\mathbf{R}}\) 。
\end{enumerate}

同样地，当 \((x,y)\) 趋于 \((0,+\infty)\) （同时保持在 \(R\times R\) 中）时，xy 的聚点所成的集是 \(\overline{R}\) 。

\begin{enumerate}
\def\labelenumi{\arabic{enumi})}
\setcounter{enumi}{5}
\item
试证：每个有理函数 \(\mathbf{P}(x)/\mathbf{Q}(x)\) （系数在 R 中），只要它在所有点 x （ R 的、使 \(\mathbf{Q}(x)\neq0\) 的）处有定义，就可连续延拓到点 \(+\infty\) 与 \(-\infty\) （取值于 \(\overline{R}\) ）。
\item
设 X 是一个全序集。
\end{enumerate}

\begin{enumerate}
\def\labelenumi{\alph{enumi})}
\item
设 \(X_{1}\) 是 X 的完备化（集合论，第三章，§ 1，习题 15）。\(X_{1}\) 是一个全序集，其元素是 X 的这样的子集 A：(i) 若 \(x \in A\) 且 \(y \leqslant x\) ，则 \(y \in A\) ；(ii) 若 A 在 X 中有上确界，则该确界属于 A{[}这等价于说 A 在 \(\mathcal{G}_{0}(X)\) 中是闭的{]}；\(X_{1}\) 的元素按包含关系排序。\(X_{1}\) 关于拓扑 \(\mathcal{G}_{0}(X_{1})\) 是紧的{[}§ 2，习题 6 a){]}。映射 \(x \to ]\leftarrow,x]\) （ X 到 \(X_{1}\) 中的）是严格保序的，并且若把 X 等同于它在这个映射下的像，则 X 在 \(X_{1}\) 中稠密，且 \(\mathcal{G}_{0}(X)\) 由 \(\mathcal{G}_{0}(X_{1})\) 诱导。\(X_{1}\) 连通当且仅当 X 无间隙（§ 2，习题 7）。
\item
对每个势 c，试给出一个全序集 X 的例子，它在拓扑 \(\mathfrak{G}_{0}(\mathbf{X})\) 下连通，并且使得 X 的任何一点的任何基本邻域系都具有势 \(\geqslant c\) {[}利用 a) 与 §1 习题 4 b) 的方法{]}。
\item
用 \(a\) 的记号，令 \(\mathbf{X}_{2}\) 为字典积 \(\mathbf{X}_{1} \times \left\{-\mathbf{i}, 0, \mathbf{i}\right\}\) 的这样的子集：它是由下列诸点组成的集的余集：(i) 点 \((x, 0)\) ，其中 \(x \notin \mathbf{X}\) ；(ii) 点 \((x, 1)\) ，其中 \(x \in \mathbf{X}\) 且所有 \(y > x\) （ \(\mathbf{X}\) 中的）所成之集有最小元素；(iii) 点 \((x, -1)\) ，其中 \(x \in \mathbf{X}\) 且所有 \(y < x\) （ \(\mathbf{X}\) 中的）所成之集有最大元素。试证：\(\mathbf{X}_{2}\) 赋予拓扑 \(\mathfrak{C}_{0}(\mathbf{X}_{2})\) 后是紧且全不连通的，\(\mathbf{X}\) 可借助严格保序映射 \(x \to (x, 0)\) 等同于 \(\mathbf{X}_{2}\) 的一个稠密子集，并且 \(\mathfrak{C}_{0}(\mathbf{X}_{2})\) 在 \(\mathbf{X}\) 上诱导的拓扑是离散的。
\item
设 \(X'\) 是一个包含 X 的全序集，它在 X 上诱导给定的序结构，在拓扑 \(\mathcal{C}_{0}(\mathbf{X}')\) 下是紧的，并且 X 关于这个拓扑在 \(X'\) 中稠密。试证：存在一个连续保序的满射映射 \(f: X' \to X_{1}\) 与一个连续保序的满射映射 \(g: X_{2} \to X'\) ，两者都在 X 上诱导恒等映射。
\end{enumerate}

